\section{Writing Style}
\label{sec:style}
The primary goal of technical writing is to communicate ideas precisely
and efficiently. In a CS or engineering report, clear language is more
valuable than sophisticated-sounding language. The following guidelines
cover the main aspects of technical style: clarity, tense, voice, precision,
terminology, and paragraph structure.

\subsection{Clarity over cleverness}
Prefer the plain word to the fancy one, and the short sentence to the
long one. ``We used'' beats ``We made utilization of''. ``The algorithm
runs in $O(n\log n)$ time'' beats ``The algorithm exhibits logarithmic
scaling behaviour with respect to input cardinality''. If a sentence
takes two readings to parse, split it.

\subsection{Tense}
As a simple default:
\begin{itemize}[nosep]
      \item \textbf{Past tense} for what you did and what you observed:
            ``We evaluated three schedulers'', ``The cache hit rate dropped
            to 62\,\%''.
      \item \textbf{Present tense} for facts that remain true and for
            describing the content of your document itself: ``\Cref{fig:latency}
            shows the results'', ``TCP guarantees in-order delivery''.
\end{itemize}

\subsection{Passive vs.\ active voice, and first person}
Both active and passive voice are acceptable in engineering writing;
the goal is clarity, not a blanket rule. Active voice is usually more
direct: ``We measured latency under three loads'' is easier to parse than
``Latency was measured under three loads by us''. Passive voice is useful
when the actor is irrelevant or obvious: ``The packets were captured
using \texttt{tcpdump}'' is fine because \emph{who} ran the tool matters
less than \emph{how} the data was captured. Using ``we'' for actions your
team actually performed is standard in modern CS writing --- check your
course or venue's convention if unsure.

\begin{pitfallbox}
      Do not mix an impersonal, passive style in the results with a chatty
      ``I think'' style in the discussion. Pick a register and keep it
      consistent across the whole report.
\end{pitfallbox}

\subsection{Precision and hedging}
Say exactly what you mean and no more. ``The system is fast'' is not a
claim you can defend; ``the system processes 12{,}400 requests/s at the
95th-percentile latency of 8\,ms on the hardware described in the
experimental setup section'' is. Conversely, do not overclaim: if you tested on
one dataset, say so, and do not generalise to ``all workloads''. Words
like ``clearly'', ``obviously'', and ``trivially'' rarely help the
reader and often hide a step you skipped.

\subsection{Consistency}
Pick one term for each concept and stick to it. If your system is a
``scheduler'' in \cref{sec:structure}, do not call it a ``dispatcher'' in
the evaluation unless you have explicitly defined them as synonyms.
The same applies to notation: if $n$ is the number of nodes in
\cref{sec:algorithms}, it must still be the number of nodes on the last
page.

\subsection{Paragraph and section structure}
One idea per paragraph, and put the main point in the first sentence
(a ``topic sentence''). Readers --- and graders --- skim; a report whose
first sentence of every paragraph forms a coherent summary is a
well-structured report. Use subsections to help navigation, but avoid
sections with a single sentence and no subsections around them; if a
subsection is that short, it is probably a paragraph, not a subsection.

\begin{tipbox}
      Read your own paragraph's first sentence out loud, in order, skipping
      everything else. If that summary doesn't make sense as an outline of
      your argument, restructure the paragraphs.
\end{tipbox}
